\definicion{Minificación}{Proceso de eliminar espacios y acortar nombres en el código para reducir el tamaño del archivo y mejorar la velocidad de carga, no para ocultar la lógica.}
\subtituloA{Principio de seguridad web}

\lista{
  \textbf{Exposición del código:} El navegador descarga el código completo, haciéndolo visible a cualquiera. |
  \textbf{Validación de seguridad:} Debe realizarse obligatoriamente en el servidor; la validación en el cliente es solo para comodidad del usuario. |
  \textbf{Claves sensibles:} Nunca deben incluirse en el código del cliente (ej. claves de API, contraseñas).
}
\nota{Para el examen}{
¿Por qué una clave de API no debe estar en el código cliente? Porque el navegador descarga el código completo, exponiendo la clave públicamente.
}
